\begin{tikzpicture}
\path [use as bounding box] (-1, -1) rectangle (35, 35);
    \coordinate (I) at (10.0, 18.0);
    \coordinate (M) at (29.0, 26.0);
    \coordinate (S) at (29.0, 20.0);
    \coordinate (J) at (35.0, 20.0);
    \filldraw[fill=violet, draw=black, thick] (J) -- (I) -- (M) -- (S) -- cycle;    
    \node at (10.0, 17.8) [font=\small, color=black] {I};
    \node at (29.0, 26.2) [font=\small, color=black] {M};
    \node at (29.0, 20.2) [font=\small, color=black] {S};
    \node at (35.0, 20.2) [font=\small, color=black] {J};
\end{tikzpicture}

\title{
Information on liusmedel som inte är färdigförpackade
}

\title{
VIKTIGT OM ALLERGENER
}
\section*{Vad är allergener?}
Allergener är ämnen, vanligtvis proteiner; som orsakar en allergisk reaktion hos den som är allergisk. Symtomen kan variera från kåda och magont till mer allvarliga symtom, till exempel anafylaktisk chock.
I västra fall kan symtomen bli så allvarliga att de leder till döden. Det är därför mycket viktigt att säkerställa att informationen, som ges om allergena ingredienser från allergenilistan, är korrekt. Att dokumentera vad en maträtt innehåller är viktigt för att kunna ge rätt information till konsumenten. Tänk också på att gå genom hela ingredienslistan för varje sammanst att ingrediens, till exempel pasta och majnansö, dich faga alltid leverantören en extra gång om du är osäker:
Ska framhaväs i märkningen
Från och med den 13 december 2014 är det obligatoriskt för livsmedelsföretagare att i märkningen framhäva ingredienser från allergenlistan. Framhänningen kan exempelvis göras med fet text. VERSALER eller kursiv text. Livsmedel, som har märkts före den 13 december 2014, får säljas tills lagret är slut. Därför kommer det att finnas livsmedel som är märkta enligt de gamba reglerna under en lång tid framöver. Det kommer alltså även efter den 13 december finnas livsmedelsförpackningar där ingredienserna från allergenlistan inte är framhävda.
\section*{Yarn oga med hanteringen}
För känsliga allergier och glutenintoleranta kan även mycket små mängder till exempel en droppe mjölk, en nötflsia eller några brödsmulor orsaka allergisk reaktion. Därför är det viktig att ha rutiner som förhindrar att ingredienser överförs från ett livsmedel till ett annat.
de fall konsumentens eventuella allergier har utretts i förväg ska de känna sig helt säkra på att maten är särskilt anpassad till deras behov utan att fägå. I de fall där man känner till vissa, men inte alla magtästers allgerjer och överkänslighet, behöver uppmaningen om allergeninformationen vara tydligt rikdat till de som inte får specialkost. Specialkost som är särskilt anpassad för en allergisk eller överkänslig magtäst och anges vara "fri från" ett eller flera allergen kräver precis som tidgare särskilda rutiner för tillredning och servering.
geninformationen finns att få. Fortfarande ska samtlig obligatorisk information alltid kunna ges på begäran. Se även avsnitter "Hur informationen ska lämnans".
\section*{Information mellan företag}
Om du leverarer livsmedel som inte är färdigförpackade till andra företag är du skyldig att ge nödvändig livssmedelsinformation till det mottagande företaget. Med nödvändig information menas all den information det mottagande företaget behöver för att si tun kurna ge rätt information till konsumenten. För att vetta vilka uppgifter du ska lämna vidare behöver du få information om hur livsmedlet ska säljas, dvs som oförpackat, förpackat på försälningsstället på konsumentens begäran eller färdigförpackat för direkt försäljning.
Se även avsnitter om vilken information som ska lämnas till konsumenten.
\section*{Tabell I. Allergenlistan}
I. Spannnäl som innehåller gluten, dvs. vette, räg, kom, havre, spelt, khorasanvete eller hybridiserade sorter därav och produkter därav
2. Kräftdjur och produkter därav
3. Ägg och produkter därav
4. Fisk och produkter därav
5. Jorndötter och produkter framställda därav
6. Sojäbörn och produkter därav
7. Mjölk och produkter därav (lnkusive laktos)
8. Nötter såsom mandel, haselnöt, valnöt, cashewnöt, pekanönt, paranönt; pistsandmel, makadamianöt/Queenslandsnöt, öppet
9. Selleri och produkter framställda därav
10. Senap och produkter framställda därav
11. Sesamfrön och produkter framställda därav
12. Svædlöxoid och sulfit som överstiger en viss halt
13. Lupin och produkter framställda därav
14. Blötdjur och produkter framställda därav
Undantagen ingår inte i denna sammanfattning, jämför med bilaga
III i förordning (EU) nr II 69/2011.

Information som ska ges om konsumenten frågar:
- Beteckning - namn som beskriver vad det är för slags livsmedel
- Ingredienser - informationen behöver dock inte ges i bestämd ordningsföljd
- Nettokvantitet
- Bäst före-dag eller sista förbrukningsdag
- Särskilda villkor för förvaring och/eller användning
- Kontaktuppgiff - till företaget som tillgagat och/eller förpackat maten
- Uppgift om ursprungsland eller härkomstplats (för vissa livsmedel)
- Bruksanvisning (om det behövs)
- Alkoholhalt uttryckt i volym (för drycker som innehåller mer än 1,2 volmprocent alkohol)
Hur informationen ska lämnas
Informationen om livsmedel, som inte är färdigföpragcade, får antingen ges skriftligt eller multigt. Om du väljer att ge information om allergener multigt måste det alltid finnas en särskild upplysning till konsumenten om att informationen finns att få. Det kan exempelvis finnas en skylt vid serveringen där det står "Ar du allergisk? Du kan alltid fråga oss om innehållet i maten". Det måste vara möjligt att kontrollera att informationen som ges är korrekt. Därför kan det vara nödvändigt att den som lagar maten dokumentarer vilka ingredienser som tillsätts.
Ibland kan en matgästs eventuella allergier ha utretts och noterats i förväg. Detta är vanligt i exempelvis skolor, förskolor eller på sjukhus. Till dessa magtäster krävs ingen särskild upplysning om att allergeninformationen finns att få.
Allergeninformationen och övrig livsmedelsinformation får även lämnas skriftligt, exempelvis enligt något av följande alternativ:
- På en skylt i närheten av livsmedlet
- I menyer
- I en skrift som följer med livsmedlet
- På en läplatta eller datorskärm
Om den enda allergen, som ingår i ett livsmedel, utgör en del av livsmedellets beteckning, till exempel sellerisoppa, är det tillräckligt att ange beteckningen.
\section*{Distansförsäljning}
Livsmedel som säljs via distans, exempelvis via internet eller telefon, kan till exempel vara färdigförpackade - som en burk crème fraiche i en middagskasse - eller förpackade på försäljningsstället på konsumentens begäran - till exempel tomater som förpackas och läggs ner i samma middagskasse varpå den levereras hem till konsumentens bostad.
Vid distansförsäljning gäller att konsumenten ska få samma information innan köpet som är möjligt att få vid köp i exempelvis en livsmedelsbutik. Undantag finns för information om åst före-dag eller sista förbrukningsdag, vilket måste kunna lämnas först när livsmedelt levereras. Information om bast före-dag eller sista förbrukningsdag ska för livsmedel, som är förpackade på försäljningsstället på konsumentens begäran, kunna lämnas om kunden frågar. Detsamma gäller för övrigg upspifter som är obligatoriska för denna kategori av livsmedel, ut information om allergener. Information om allergener ska alltid lämnas innan köpet avslutas. Se avsnittet "Vilken information ska lämnas till konsumenten?"
Hur informationen kan lämnas innan köpet
Det är möjligt att informera konsumenten innan köpet på olika sätt. Några exempel är via en webbsida eller genom att hänvisa till ett telefonnummer till samnans med tillgängliga telefontidet. Konsumenten ska kunna få informationen utan att dekot något extra, varför det exempelvis inte går att hänvisa till ett betalnumen.
Att lämna information vid leverans av varan
Vid leverans av varan han information om allergener, som alltid ska ges, och övriga obligatoriska uppgifter, som ska kunna lämnas på förfrågan, lämnas skriftligt eller multigt. Om informationen ges multigt krävs en särskild upplysning om att allergeninformationen finns att få. Om den som levererar maten inte kan informera om allergener och övriga obligatoriska uppgifter kan det vara lämpligt att informationen finns skriftligt eller, om det är praktiskt möjligt, att leverantören kan kontaka dan som tillagat maten. För matlädor som levereras till personer i deras hem under en längre period är det vanligt att information om kundens eventuella allergier tas om hand vid beställningen av tjänsten. I sådana fall krävs ingen särskild upplysning om att aller-

\title{
L\&S MER
}
\section*{Vagledningar m.m.}
P\&L Lismedselverkets webblts www.lismedselverket.se Informationsforordningen
Om allergi och annan överkänslighet
P\&E U-kommissionens webblts
Fragor och var som tillämpningen av förordning (EU) n 1169/2011 om tillhandallhändande av
P\&L förbruksverkets till konsumterna
P\&Jörbruksverkets webblts
Om handelsnomer för frut och grönsaker
Banschriktetniger
Lismedselindustrins och dagligvarhundelns branschiktlinjer för allergi och annan överkänslighet - Hanteering och märkning av lismedsel
Lagsführung om livmedelsinformation
- Förordning (EU) n 1169/2011 om tillhandabhållande av
lismedelsinformation till konsumenterna
- Förordning (EU) n 1379/2013 om konsumentinformation och
märkning av fikerpröudrukter; vattenbruksvordkuter; siggår's och alger
vattenbruksvordkuter
- Förordning (EU) n 1337/2013 om angivande av ursprungland eller
härkomstplats för färst, ktyt eller fryst kött av svin, får get och föäferf
- Förordning (EC) n 1760/2000 om upprättande av ett system för
- Förordning (EC) n 1760/2000 omm upprättande av notvärkut särmt mätkning av nöktött och nöktöttsprudter
- Lismedselverkrets föreskrifter (LVSVS 2014-4) om lvismedelsinforma-
nöktittsprudter
- Lismedselverkrets föreskrifter om handelsbeteckningar på fiskerioch vattenbruksvordkuter SLVFS 2001:37, senaste ändring av bilgan LVSVS 2014:23
- Mer lagstiftning som bland annat rör lismedelsinformation om olika lag: lismedel, hittar du på Lismedelverkrets webblts

\title{
Exempel på hur informationen kan ges skriftligen
}

\title{
Att informera konsumenterna om livsmedel som inte är fördigförpackade
}
\author{
Vilken information ska lämnas \\ till konsumenten? \\ Oavsett hur livsmedlet säljs, dvs. om det säljs ofö \\ packat, färdigförpackat för direkt försäljning eller för \\ packat på försäljningsstället på konsumentens begåran, \\ ska information om ingående allergena ingredienser och \\ processhälpmedel alltid ges. De allergena ämnen eller \\ produkter dåvar som alltid ska anges finns i informa \\ tionsförordningens bilaga II (den sk. allergenlistan). \\ Allergenlistan presenteras i tabell 1. Exempel på \\ dukter därav är laktos från mjölk, vassle från mjölk, \\ lecitin från soja och fiskrom från fisk. \\ Det är dessutom obligatoriskt att informera om \\ livsmedlets beteckning (namn som beskriver vad det \\ är för livsmedel) om konsumenten frågar efter detta. \\ Observera att för fiskeri- och vattenbruksprodukter är \\ det obligatoriskt att ange handelsbeteckning och artens \\ vetenskapliga namn (latinsk namn), till exempel på \\ anslagstavlor och planscher. \\ För livsmedel som förpackas på försäljningsstället på \\ konsumentens begåran eller är färdigförpackade för å \\ rekt försäljning ska dessutom ett ytterligare antal upp- \\ gifter kunna lämnas om konsumenten frågar (se även \\ broschyren "märkning av färdigförpackade livsmedel" \\ för mer information om respektive uppgift nedan): \\ Information som alltid ska ges: \\ - Allergener - ingredienser från allergenlistan
}

I den här broschyren finns information för dig som säjer eller serverar livsmedel som inte är färdigförpackade. Informationen kan exempelvis vara användbar för dig som arbetar i restaurang, på kafé, i ett offentligt kök eller i en manuell disk i en livsmedelsbutik. Det är viktigt att känna till att den här informationen inte är heltäckande, utan att det kan finnas andra regler som du måste ta hänsyn till beroende på vilket livsmedel det handlar om.
Du kan därför behöva komplettera dina kunskaper genom att läsa mer om andra regler. Läs alltid informationen i denna broschyr tillsammans med lagstiftningen.
\section*{Du ansvarar för att maten är säker} och att rätt information ges
Alla som serverar eller säjer livsmedel har ett ansvar för att maten är säker att åta. Livsmedel får inte medföra risk för människors hälsa.
Korrekt angiven livsmedelsinformationen kan vara livsviktig för konsumenten. Därför finns nu skärpta krav i lagstiftningen om allergeninformation, det vill säga information om maten innehåller de ingredienser som det är vanligast att vara allergisk eller överkänslig mot. Dessa ämnen listas i lagstiftningen (se nedan) samt i tabell 1.
\section*{Ta reda på vilka regler som gäller}
Reglerna är till för att skydda konsumenten Det finns två huvudsyften med livsmedelslagstiftningen:
- att skydda konsumenten från livsmedel som är otjänliga eller skadliga för hälsan och
- att ge konsumenten tillräcklig information om ett livsmedel för att han eller hon ska kunna göra egna och säkra val. Livsmedelsinformation, märkning, rekklam och presentation av livsmedel får inte vilseleda konsumenten.
Reglerna ska också underlätta för den fria rörligheten på den inre EU-marknaden.
\section*{Det här gäller för alla livsmedelsföretag:}
- Ansvar - Du som företagare ansvarar för de livsmedel som du tillverkar, transporterar, lagrar eller säjer.
- Spärbarhet - Det ska gå snabbt att identifiera en leverantör och ett parti.
- Öppenhet - Du ska omedelbart informera behöriga myndigheter och konsumenter om det finns anledning att misstänka att ett livsmedel inte är säkert.
- Förebyggande - Du ska regelbundet se över kritiska punkter i din verksamhet
- Kontroll - Du ska se till att verksamheten uppfyller målen i lagstiftningen.

\title{
Livsmedelsinformation till konsumenterna
}
Reglerna om livsmedelsinformation, som bland annat omfattar märkning på livsmedelsförpackningar, är till för att konsumenten lättare ska kunna göra medvetna val av livsmedel.
\section*{Var finns reglerna?}
Reglerna om livsmedelsinformation finns i förordning (EU) nr 1169/2011 (den så kallade informationsförordningen). Yttertligare regler om information som ska ges för livsmedel som inte är färdigförpackade finns i Livsmedelsverkets föreskrifter LIVSFS 2014:4 om livsmedelsinformation. Det finns även fler regler som gäller för oförpackade livsmedel av vist slag, exempelvis för fisk, kalvkört, nötkött, frukt och grönsaker etc.
Vad är ett färdigförpackat livsmedel?
Ett färdigförpackat livsmedel är ett livsmedel som placerades i en förpackning innan det nådde konsumenten. Förpackningen ska täcka livsmedlet helt eller delvis. Det ska inte gå att ändra innehållet utan att förpackningen öppnas eller ändras.
För färdigförpackade livsmedel körss livsmedelsinformation i form av märkning på livsmedelsförpackningen. Du kan läsa mer i broschyren "Märkning av färdigförpackade livsmedel".

\title{
LIVSMEDELS
} VERKET

\section*{Livsmedel som inte är färdigförpackade}
Livsmedel som säljs och inte är färdigförpackade delas in i tre kategorier:
\section*{Oförpackade livsmedel}
Exempelvis mat som serveras på en restaurang eller godis för lösviktsförsäljning i en butik.
Livsmedel som förpackas på försäljningsstället på konsumentens begäran
Mat eller dryck som först är oförpackad, men som placeras i en förpackning på konsumentens begäran. Livmedel som säljs via distans kan också gå in under denna kategori, exempelvis att kunden ringer och beställer mat som levereras till kundens bostad.
Livsmedel färdigförpackade för direkt försäljning Mat eller dryck som säljs på ett sätt där kunden har direkt personlig kontakt med personal som kan ge den obligatoriska livsmedelsinformationen i en eller annan form. Se vidare under rubriken "Hur informationen ska lämnas".
För livsmedel som inte är färdigförpackade får livsmedelsinformationen lämnas på ett flexibelt sätt. Villken information som ska lämnas och hur förklaras i denna broschyr.
Livsmedelsinformationen ska vara på svenska
Huvudregeln är att livsmedelsinformationen ska vara på svenska, men det går bra att använda annat språk för vissa uppgifter om skillnaden i stavning är obtydlig Exempel: ”öl” och ”öl” eller ”bär” och "bær”.
\section*{EXEMPEL PÅ LIVSMEDEL SOM INTE BETRAKTÅS SOM FÄRDIGFÖRPACKADE}
\section*{Livsmedel}
- Smörgåsar bereds och förpackas i köket på ett sjukhus för att sedan levereras till en möteslokal i annan del av sjukhuset.
- En restaurang klyer ner mat som förpackas i maltädor och säljs vid ett senare tillfälle på samma restaurang.
- Pizza som en kund ringer och beställer för hemlevarens. Leveransen kan ske via annat företag.
- Varma eller kalla mallädör som levereras från ett kommunalt kök direkt till personer i deras hem. Leveransen kan ske via annat företag.
- Ost delas och packeteras i förväg bokom en manuell disk i en butik varpå den överlämnas till kunden.
- Sill som placeras i en förpackning i kundens návarro och lämnas över en manuell disk.
- Kunden lagger själv ner sin mat i en låda från en buffé i en butik.
\section*{Kategori}
Förpackat på försäljningsstället på konsumentens begäran.
Förpackat på försäljningsstället på konsumentens begäran.
Förpackat på försäljningsstället på konsumentens begräran.
Förpackat försäljningsstället på konsumentens begäran.
Förpackat livsmedel

\title{
Information om maträtter och allergener på restaurang
}